\section{Transacciones del Subsistema de Usuarios (Fernando)}


\subsection{Transacción: Alta de usuario}

\textbf{Descripción:} Inserta un nuevo usuario en la tabla \texttt{USUARIO}. Si la inserción es válida se confirma con \texttt{COMMIT}. Ante un error (por ejemplo, duplicidad de email o nombre), se revierte la operación con \texttt{ROLLBACK}.

\begin{lstlisting}[language=SQL, caption={Transacción de alta de usuario}]
INSERT INTO USUARIO (NOMBREUSUARIO, EMAIL, CONTRASENIA, ESADMIN, BORRADO, FECHAMODIFICACION)
VALUES ('FERNANDO', 'fernando@email.com', 'contrasenia_hasheada', 0, 'N', SYSDATE);

COMMIT;
\end{lstlisting}

\subsection{Transacción: Creación de relación de amistad}

\textbf{Descripción:} Inserta una nueva relación en \texttt{AMISTAD}. Si se produce un error, se revierte la operación.

\begin{lstlisting}[language=SQL, caption={Transacción de creación de amistad}]
INSERT INTO AMISTAD (IDUSUARIO1, IDUSUARIO2, FECHAAMISTAD)
VALUES (10, 20, SYSDATE);

COMMIT;
\end{lstlisting}

\subsection{Transacción: Bloqueo de usuario y eliminación de amistad previa}

\textbf{Descripción:} Se ejecutan dos operaciones: eliminar posible amistad previa e insertar el bloqueo. Ambas deben realizarse dentro de la misma transacción. Solo se confirma si ambas se ejecutan sin errores.

\begin{lstlisting}[language=SQL, caption={Transacción de bloqueo con consistencia relacional}]
DELETE FROM AMISTAD
WHERE (IDUSUARIO1 = 10 AND IDUSUARIO2 = 20)
   OR (IDUSUARIO1 = 20 AND IDUSUARIO2 = 10);

INSERT INTO BLOQUEO (IDUSUARIO1, IDUSUARIO2, FECHABLOQUEO)
VALUES (10, 20, SYSDATE);

COMMIT;
\end{lstlisting}

\subsection{Transacción: Baja lógica (borrado lógico) de usuario}

\textbf{Descripción:} Se actualiza el campo \texttt{BORRADO} y la \texttt{FECHAELIMINACION}. La operación se confirma con \texttt{COMMIT}.

\begin{lstlisting}[language=SQL, caption={Transacción de baja lógica de usuario}]
UPDATE USUARIO
SET BORRADO = 'Y',
    FECHAELIMINACION = SYSDATE,
    FECHAMODIFICACION = SYSDATE
WHERE IDUSUARIO = 10;

COMMIT;
\end{lstlisting}
